For the literal Construction3 on the classical unital UH(q), the design
parameters are 2-(q³+1,q+1,1), and each point is on q² blocks.
Therefore:

\begin{itemize}
\tightlist
\item
  point objects: q³+1;
\item
  marked flag objects: q²(q³+1);
\item
  block objects: q²(q²−q+1);
\item
  panel sizes, in Figure2's point/marked/block order: q, q+1, q²−1;
\item
  local orders printed under those nodes should be q−1, q, q²−2.
\end{itemize}

These are direct counts of the actual source incidences. They do not
depend on the claimed all-q connectivity or on a choice of acting group.

\subsubsection{Panel sizes}\label{panel-sizes}

Three printed rows have smaller marked/block panel sizes:

{\def\LTcaptype{none} % do not increment counter
\begin{longtable}[]{@{}rll@{}}
\toprule\noalign{}
q & Printed local orders & Literal local orders \\
\midrule\noalign{}
\endhead
\bottomrule\noalign{}
\endlastfoot
5 & 4,1,7 & 4,5,23 \\
8 & 7,2,20 & 7,8,62 \\
11 & 10,3,39 & 10,11,119 \\
\end{longtable}
}

In each of these rows, the printed marked panel size is (q+1)/3 and the
printed block panel size is (q²−1)/3. Thus the diagram cannot describe
the full literal Construction3 with the printed object set.

This factor three agrees with {[}PGU(3,q):PSU(3,q){]}=gcd(3,q+1) in
precisely these rows. A selected chamber orbit under the smaller group
can retain the transitive object sets while omitting incidences of the
complete construction. That is a possible explanation of the table, not
a verified reconstruction of the source's computational procedure. No
claim that the author used a particular algorithm is made here.

\subsubsection{Object counts at q=16}\label{object-counts-at-q16}

The printed row gives 1,028 point objects, 263,206 marked objects and
15,482 block objects. Literal unital parameters instead give:

{\def\LTcaptype{none} % do not increment counter
\begin{longtable}[]{@{}lrr@{}}
\toprule\noalign{}
Object type & Printed & Literal Construction3 \\
\midrule\noalign{}
\endhead
\bottomrule\noalign{}
\endlastfoot
Points & 1,028 & 4,097 \\
Marked flags & 263,206 & 1,048,832 \\
Blocks & 15,482 & 61,696 \\
\end{longtable}
}

The printed row is internally incompatible with its own q=16 flag
valencies: 1,028·256=263,168 and 15,482·17=263,194, neither equal to
263,206. This is a numerical correction independent of any
interpretation of the group-index issue above.

All other printed object counts for q=3,4,5,7,8,9,11,13 agree with the
literal formulas. The local orders in rows q=3,4,7,9,13,16 agree with
q−1,q,q²−2.
